\begingroup
\scriptsize
\setlength{\tabcolsep}{2pt}
\begin{longtable}{@{}p{.055\linewidth}p{.12\linewidth}p{.235\linewidth}p{.225\linewidth}p{.225\linewidth}p{.07\linewidth}@{}}
\caption{Estimand-derived identification requirements. G denotes generic benchmark validity; S denotes a security procedural-transfer specialization.}\label{tab:contract}\\
\toprule
ID & Short name & Definition & Causal interpretation licensed & False inference if absent & Map \\
\midrule
\endfirsthead
\toprule
ID & Short name & Definition & Causal interpretation licensed & False inference if absent & Map \\
\midrule
\endhead
\multicolumn{6}{@{}l}{\textbf{A. Target / outcome validity}} \\
A1 & Nontrivial target & The requested target functionality is genuinely absent before intervention and a no-op does not satisfy the task. & Licenses interpreting later task success as newly achieved rather than pre-existing or no-op success. & A passing endpoint can be mistaken for implementation even when the requested behavior was already present. & G \\
A2 & Insecure completion & There exists an implementation/state that genuinely completes the requested task while failing the focal security requirement. & Shows that task completion and focal insecurity can coexist, so the harmful endpoint is attainable. & Security failure can be conflated with failure to implement the requested task. & S \\
A3 & Secure completion & There exists an implementation/state that genuinely completes the same requested task while passing focal security. & Shows that the requested task can be completed while satisfying the same focal security property. & The design can mistake an impossible security requirement for agent failure or treatment harm. & S \\
A4 & Feature retention & The secure repair preserves the requested functionality rather than removing the feature or changing the task. & Licenses treating the secure state as a repair of the same task rather than removal of the feature. & Deleting or weakening the requested behavior can masquerade as a secure repair. & S \\
A5 & Focal witness & The focal security endpoint is executable or otherwise independently grounded and discriminates the intended unsafe/safe states. & Makes the security component of task-complete insecurity an independently grounded outcome. & A proxy or nondiscriminating test can label states safe or unsafe without measuring the focal property. & G \\
\multicolumn{6}{@{}l}{\textbf{B. Source / transfer validity}} \\
B1 & Source correctness & The source procedure is demonstrated to work in its source context. & Supports describing the donated procedure as one that worked in its original context. & A broken source episode can be mischaracterized as successful procedural knowledge. & G \\
B2 & Source focal safety & The source procedure is demonstrated to be safe with respect to the focal security property in its original source context. & Supports the claim that a focal-safe source procedure, rather than an already unsafe one, is transferred. & Target insecurity can be attributed to lost applicability when the donor was never focal-safe. & S \\
B3 & Procedural relevance & The source procedure and target task genuinely share the procedural operation whose transfer is being studied. & Connects the delivered source episode to the operation required by the target task. & Topical or lexical similarity can be mistaken for transfer of the operative procedure. & S \\
B4 & Applicability predicate & A named condition/precondition governing safe transfer is specified. & States the condition under which the relevant source procedure is expected to transfer safely. & Relevance can silently stand in for applicability, leaving the proposed failure mechanism unspecified. & S \\
B5 & Directional predicate shift & The relevant condition is established in the source and differs/fails in the target. & Grounds the intended source-to-target change by showing the named predicate true at source and false at target. & A merely asserted context difference can be mistaken for the security-relevant applicability shift. & S \\
B6 & No major co-mismatch & The intended applicability/security difference is not obviously confounded by another comparably important source-target incompatibility. & Makes the named applicability shift the salient pair-level incompatibility rather than one of several alternatives. & A second framework, API, or task mismatch can explain failure attributed to the focal predicate. & S \\
\multicolumn{6}{@{}l}{\textbf{C. Treatment validity}} \\
C1 & Delivery fidelity & The treatment actually delivers the intended source memory/skill in a defined and reproducible form. & Defines a reproducible assigned intervention and verifies delivery of the intended source episode. & The analyzed treatment can differ from the memory or skill the causal claim names. & G \\
C2 & Resource balance & The relevant comparison controls, measures, or explicitly treats context, token, action, tool, and runtime differences introduced by memory/skill delivery. & Separates or explicitly defines semantic content effects relative to context, action, tool, and runtime burden. & Resource displacement caused by delivery can be attributed to procedural semantics. & S \\
C3 & Semantic control & The evaluation includes a comparison capable of separating the intended memory/skill content effect from generic additional-context/skill-loading effects where such separation is part of the stated claim. & Provides the comparator required to isolate relevant memory content from generic extra-context effects. & Any effect of loading additional material can be called an effect of the relevant procedure. & S \\
C4 & Outcome-blind source lock & The source/skill used for the evaluated target is not chosen or changed after observing the target's evaluated outcome. & Prevents target outcomes from influencing which donor procedure becomes the treatment. & Outcome-guided treatment selection can be mistaken for prospective causal evidence or frequency. & S \\
\multicolumn{6}{@{}l}{\textbf{D. Process / analysis validity}} \\
D1 & Delivery / uptake separation & The study distinguishes treatment delivery from evidence that the agent actually used/engaged with the memory or skill when mechanism claims require that distinction. & Supports mechanism claims by distinguishing assigned and delivered memory from behavioral engagement with it. & Delivery alone can be reported as procedural use, or post-treatment uptake can be treated as baseline eligibility. & S \\
D2 & Joint task-security outcome & The analysis distinguishes task failure from task-complete insecurity rather than interpreting security conditional on a post-treatment success state in a way that changes the estimand. & Identifies task-complete insecurity directly without conditioning security on a treatment-responsive success variable. & Separate marginals or selected successful runs can answer a different causal question. & S \\
D3 & Valid units and replication & Reported repetitions/seeds/tasks are interpreted using defensible experimental units rather than treating effectively deterministic duplicates as independent evidence. & Aligns uncertainty and evidence strength with distinct tasks or genuinely varying executions. & Deterministic duplicates can be counted as independent evidence and yield spurious precision. & G \\
\bottomrule
\end{longtable}
\endgroup
